\documentclass[10pt]{mypackage}

\usepackage[uprightgreeks,light,noDcommand]{kpfonts}
\renewcommand{\coloneq}{\coloneqq}
\renewcommand{\eqcolon}{\eqqcolon}
%\renewcommand{\mathbb}{\mathds}
\usepackage{homework}
%\usepackage{notes}

%\usepackage[ backend=bibtex, style = alphabetic, sorting=ynt ]{biblatex}
%\addbibresource{  }

\usepackage{parskip}

\fancyhf{}
\fancyhead[R]{Avinash Iyer}
\fancyhead[L]{Algebra II: Homework 11}
\fancyfoot[C]{\thepage}

\setcounter{secnumdepth}{0}

\begin{document}
\RaggedRight
\begin{problem}[Problem 1]
  Let $p$ be a prime, $n$ a positive integer, and let $\Phi_n = x^{p^{n}}-x$ be an element of $\F_p[x]$. Prove that $\Phi_n$ is equal to the product of all the monic irreducible polynomials in $\F_p[x]$ whose degrees divide $n$, with each polynomial appearing with multiplicity one.
\end{problem}
\begin{solution}
  Fix an algebraic closure $ \overline{\F_p} $, and view any polynomial as $\F_p[x]$ as splitting into linear factors with roots in $ \overline{\F_p} $.

  The splitting field of $\Phi_n(x)$ is then given by $\F_{p^{n}}$, and we may write
  \begin{align*}
    \Phi_n(x) &= \prod_{\alpha\in \F_{p^{n}}} \left( x-\alpha \right),
  \end{align*}
  observing that each factor $\left( x-\alpha \right)$ is distinct. Any irreducible factor $f_i$ of $\Phi_n$ then defines a subfield of the splitting field of $\Phi_n(x)$, given by $\F_{p^{m}}\subseteq \F_{p^{n}}$ for some $n$. In particular, we have $m | n$, so that $\deg\left( f_i \right) | n$. Since $f_i$ divides $\Phi_n(x)$, it follows that over $ \overline{\F_p} $, we have
  \begin{align*}
    f_i(x) &= \prod_{\alpha\in \F_{p^{m}}} \left( x-\alpha \right),
  \end{align*}
  and since $\F_{p^{m}}\subseteq \F_{p^{n}}\subseteq \overline{\F_p}$ is unique, it follows that $f_i$ is unique.

  Additionally, if $f_j(x)$ is any (monic) irreducible polynomial in $\F_p[x]$ of degree $m | n$, we have that $f_j(x)$ defines a field extension $F\subseteq \F_{p^{n}}\subseteq \overline{\F_p}$ with $\left[ F:\F_p \right] = m$, whence $F = \F_{p^{m}}$. Since $\F_{p^{m}}\subseteq \F_{p^{n}}$, it follows that every root of $f_j$ in $ \overline{\F_p} $ is a root of $ \Phi_n(x) $, whence $f_j | \Phi_n(x)$, and $f_j$ is unique by the earlier discussion. Therefore, we have
  \begin{align*}
    \Phi_n(x) &= \prod_{m | n} f_m(x),
  \end{align*}
  where $\deg\left( f_m \right) | n$ and $f_m(x)$ is the unique monic irreducible factor of degree $m$ for $\Phi_n(x)$.
\end{solution}
\begin{problem}[Problem 2]
  Let $K\subseteq \C$ be the splitting field of $f(x) = x^4-2$ over $\Q$.
  \begin{enumerate}[(a)]
    \item Choose an order on the set of roots of $x^{4}-2$ and describe the associated embedding of $\gal\left( K/\Q \right)$ into $S_4$.
    \item Describe all subgroups of $\gal\left( K/\Q \right)$ and the corresponding subfields of $K$.
  \end{enumerate}
\end{problem}
\begin{solution}\hfill
  \begin{enumerate}[(a)]
    \item We observe that $K = \Q\left( \sqrt[4]{2},i \right)$ has degree $8$ over $\Q$, with basis given by
      \begin{align*}
        A &= \set{1,\sqrt[4]{2},\sqrt[4]{4},\sqrt[4]{8},i,i\sqrt[4]{2},i\sqrt[4]{4},i\sqrt[4]{8}}.
      \end{align*}
      Applying a $\Q$-automorphism to any element $k\in K$ written in its basis gives
      \begin{align*}
        \sigma\left( k \right) &= k_1 + k_2\sigma\left( \sqrt[4]{2} \right) + k_3\left( \sigma\left( \sqrt[4]{2} \right) \right)^2 + k_4\left( \sigma\left( \sqrt[4]{2} \right) \right)^{3}\\
                               &+ k_4\sigma\left( i \right) + k_5\sigma\left( i \right)\sigma\left( \sqrt[4]{2} \right) + k_6 \sigma\left( i \right)\left( \sigma\left( \sqrt[4]{2} \right) \right)^2 + k_7 \sigma\left( i \right)\left( \sigma\left( \sqrt[4]{2} \right) \right)^3.
      \end{align*}
      We have that $\gal\left( K/\Q \right)$ acts transitively on the roots of $f(x)$, which are given by $\set{\sqrt[4]{2},i\sqrt[4]{2},-\sqrt[4]{2},-i\sqrt[4]{2}}$, so we consider the automorphisms defined by $\sqrt[4]{2}\xmapsto{\sigma}i\sqrt[4]{2}$ and $i\xmapsto{\tau}-i$, where $\tau\left( \sqrt[4]{2} \right) = \sqrt[4]{2}$. Observe that under this ordering, we have representations of $\sigma = \left( 1,2,3,4 \right)$ and $\tau = \left( 2,4 \right)$ as cycle decompositions in $S_4$.
    \item The subgroups of $D_4 = \left\langle \sigma,\tau|\sigma^4,\tau^2,\tau\sigma\tau\sigma \right\rangle$, where $\sigma$ and $\tau$ are defined as above, are given by
      \begin{align*}
        H_0 &= \set{e}\\
        H_1 &= \left\langle \sigma \right\rangle\\
        H_2 &= \left\langle \sigma^2 \right\rangle\\
        H_3 &= \left\langle \tau \right\rangle\\
        H_4 &= \left\langle \sigma\tau \right\rangle\\
        H_5 &= \left\langle \sigma^2\tau \right\rangle\\
        H_6 &= \left\langle \sigma^3\tau \right\rangle\\
        H_7 &= \left\langle \sigma^2,\tau \right\rangle\\
        H_8 &= \left\langle \sigma^2,\sigma\tau \right\rangle\\
        G &= \left\langle \sigma,\tau \right\rangle.
      \end{align*}
      Now, we compute fixed fields.
      \begin{itemize}
        \item First, we observe that $K^{H_0} = K$ almost by definition.
        \item To find $K^{H_1}$, we observe the action of $\sigma$ on the basis $A$ is given by
          \begin{align*}
            \sigma &= \begin{cases}
              1\mapsto 1\\
              \sqrt[4]{2} \mapsto i\sqrt[4]{2}\\
              \sqrt[4]{4} \mapsto -\sqrt[4]{4}\\
              \sqrt[4]{8} \mapsto -i\sqrt[4]{8}\\
              i \mapsto i\\
              i\sqrt[4]{2} \mapsto -\sqrt[4]{2}\\
              i\sqrt[4]{4} \mapsto -i\sqrt[4]{4}\\
              i\sqrt[4]{8} \mapsto \sqrt[4]{8}.
            \end{cases}
          \end{align*}
          Therefore, $\Q\left( i \right)\subseteq K^{H_1}$. Since $\left[ G:H_1 \right] = 2 = \left[ \Q\left( i \right):\Q \right]$, it follows that $K^{H_1} = \Q\left( i \right)$.
        \item For $K^{H_2}$, the action of $\sigma^2$ on the basis $A$ is given by
          \begin{align*}
            \sigma^2 &= \begin{cases}
              1\mapsto 1\\
              \sqrt[4]{2} \mapsto -\sqrt[4]{2}\\
              \sqrt[4]{4} \mapsto \sqrt[4]{4}\\
              \sqrt[4]{8} \mapsto -\sqrt[4]{8}\\
              i\mapsto i\\
              i\sqrt[4]{2}\mapsto -i\sqrt[4]{2}\\
              i\sqrt[4]{4} \mapsto i\sqrt[4]{4}\\
              i\sqrt[4]{8} \mapsto -i\sqrt[4]{8}.
            \end{cases}
          \end{align*}
          Therefore, we have $\Q\left( i,\sqrt{2} \right)\subseteq K^{H_2}$. Since $\left[ G:H_2 \right] = 4 = \left[ \Q\left( i,\sqrt{2} \right):\Q \right]$, equivalence follows.
        \item For $K^{H_3}$, the action of $\tau$ on the basis $A$ is given by
          \begin{align*}
            \tau &= \begin{cases}
              1\mapsto 1\\
              \sqrt[4]{2} \mapsto \sqrt[4]{2}\\
              \sqrt[4]{4} \mapsto \sqrt[4]{4}\\
              \sqrt[4]{8} \mapsto \sqrt[4]{8}\\
              i\mapsto -i\\
              i\sqrt[4]{2} \mapsto -i\sqrt[4]{2}\\
              i\sqrt[4]{4} \mapsto -i\sqrt[4]{4}\\
              i\sqrt[4]{8}\mapsto -i\sqrt[4]{8}.
            \end{cases}
          \end{align*}
          Therefore, $\Q\left( \sqrt[4]{2} \right)\subseteq K^{H_3}$, while since $\left[ G:H_3 \right]=4=\left[ \Q\left( \sqrt[4]{2} \right):\Q \right]$, we have $K^{H_3} = \Q\left( \sqrt[4]{2} \right)$.
        \item For $K^{H_4}$, the action of $\sigma\tau$ on the basis $A$ is given by
          \begin{align*}
            \sigma\tau &= \begin{cases}
              1\mapsto 1\\
              \sqrt[4]{2}\xmapsto{\tau}\sqrt[4]{2}\xmapsto{\sigma}i\sqrt[4]{2}\\
              \sqrt[4]{4}\xmapsto{\tau}\sqrt[4]{4}\xmapsto{\sigma}-\sqrt[4]{4}\\
              \sqrt[4]{8}\xmapsto{\tau}\sqrt[4]{8}\xmapsto{\sigma}-i\sqrt[4]{8}\\
              i\xmapsto{\tau}-i\xmapsto{\sigma}-i\\
              i\sqrt[4]{2}\xmapsto{\tau}-i\sqrt[4]{2}\xmapsto{\sigma}\sqrt[4]{2}\\
              i\sqrt[4]{4}\xmapsto{\tau}-i\sqrt[4]{4}\xmapsto{\sigma} i\sqrt[4]{4}\\
              i\sqrt[4]{8}\xmapsto{\tau}-i\sqrt[4]{8}\xmapsto{\sigma} -\sqrt[4]{8}.
            \end{cases}
          \end{align*}
          To find the fixed field in this scenario, we let $x = \sum_{j=1}^{8}a_je_j$ be an element of $\Q\left( \sqrt[4]{2},i \right)$, where
          \begin{align*}
            \set{e_1,\dots,e_8} = \set{1,\sqrt[4]{2},\dots,i\sqrt[4]{8}}
          \end{align*}
          is our preferred ordering of the basis over $\Q$. Then, we have
          \begin{align*}
            x &= a_1 e_1 + a_2 e_2 + a_3e_3 + a_4e_4 + a_5e_5 + a_6e_6 + a_7e_7 + a_8 e_8\\
            \sigma\tau(x) &= a_1 e_1 + a_6 e_2 - a_3e_3 -a_8e_4 - a_5e_5 + a_2 e_6 + a_7e_7 -a_4e_8
          \end{align*}
          For equality to hold, we must have $a_5 = -a_5$, $a_3 = -a_3$, $a_2 = a_6$, and $a_4 = -a_8$. Writing this out explicitly, we have
          \begin{align*}
            x &= a_1 + a_2 \left( \sqrt[4]{2} + i\sqrt[4]{2} \right) + a_7 \left( i\sqrt[4]{4} \right) + a_4 \left( \sqrt[4]{8}-i\sqrt[4]{8} \right)\\
              &= a_1 + a_2 \left( \sqrt[4]{2}\left( 1+i \right) \right) + \frac{1}{2}a_7 \left( \sqrt[4]{2}\left( 1+i \right) \right)^2 - \frac{1}{2} a_4\left( \sqrt[4]{2}\left( 1+i \right) \right)^{3}.
          \end{align*}
          In particular, this means that $\Q\left( \sqrt[4]{2}\left( 1+i \right) \right)\subseteq K^{H_4}$. Since $\left[ \Q\left( \sqrt[4]{2}\left( 1+i \right) \right):\Q \right] = 4 = \left[ G:H_4 \right]$, equality follows.
        \item Similarly for $H_5$, the action of $\sigma^2\tau$ on the basis $A$ is given by
          \begin{align*}
            \sigma^2\tau &= \begin{cases}
              1\mapsto 1\\
              \sqrt[4]{2}\xmapsto{\tau}\sqrt[4]{2}\xmapsto{\sigma^2}-\sqrt[4]{2}\\
              \sqrt[4]{4}\xmapsto{\tau}\sqrt[4]{4}\xmapsto{\sigma^2}\sqrt[4]{4}\\
              \sqrt[4]{8}\xmapsto{\tau}\sqrt[4]{8}\xmapsto{\sigma^2}-\sqrt[4]{8}\\
              i\xmapsto{\tau}-i\xmapsto{\sigma^2}-i\\
              i\sqrt[4]{2}\xmapsto{\tau}-i\sqrt[4]{2}\xmapsto{\sigma^2}i\sqrt[4]{2}\\
              i\sqrt[4]{4}\xmapsto{\tau}-i\sqrt[4]{4}\xmapsto{\sigma^2}-i\sqrt[4]{4}\\
              i\sqrt[4]{8}\xmapsto{\tau}-i\sqrt[4]{8}\xmapsto{\sigma^2}i\sqrt[4]{8}.
            \end{cases}
          \end{align*}
      \end{itemize}
  \end{enumerate}
\end{solution}
\begin{problem}[Problem 3]
  Let $p$ and $q$ be distinct primes with $q > p$, and let $K/F$ be a Galois extension of degree $pq$. Prove that
  \begin{enumerate}[(a)]
    \item there exists a field $L$ with $F\subseteq L\subseteq K$ and $\left[ L:F \right] = q$;
    \item there exists a unique field $M$ with $F\subseteq M\subseteq K$ and $\left[ M:F \right] = p$.
  \end{enumerate}
\end{problem}
\end{document}
